\section{快速入门}

\subsection{第一段代码}

下面是一个帮助你上手的简要介绍。
我们假设已经安装了 \tkzNamePack{tkz-euclide} 和 \tkzNamePack{tkz-elements} 这两个宏包。
请使用 \tkzEngine{lualatex} 引擎编译以下代码；你应当得到一条经过点 $A$ 和 $B$ 的直线。

\medskip
\noindent
\textbf{注意。}
\tkzNamePack{tkz-elements} 在 Lua 中完成所有几何定义与计算，
而 \tkzNamePack{tkz-euclide} 主要用于绘图。
因此推荐使用 \tkzNamePack{tkz-euclide} 的 \texttt{mini} 选项。
如果遇到编译问题，只需不加该选项加载 \tkzNamePack{tkz-euclide} 即可。


\begin{mybox}{}
\begin{minipage}{.55\textwidth}
  \begin{tkzexample}[code only]
  % !TEX TS-program = lualatex
  \documentclass{article}
  \usepackage[mini]{tkz-euclide}
  \usepackage{tkz-elements}
  \begin{document}
  \directlua{
    init_elements()
    z.A = point(0, 1)
    % or z.A = point:new(0, 1)
    z.B = point(2, 0)
  }
  \begin{tikzpicture}
   \tkzGetNodes
   \tkzDrawLine(A,B)
   \tkzDrawPoints(A,B)
   \tkzLabelPoints(A,B)
  \end{tikzpicture}
  \end{document}
  \end{tkzexample}
\end{minipage}
\begin{minipage}{.35\textwidth}
  \directlua{
    z.A = point(0, 1)
    z.B = point(2, 0)
  }
  \begin{tikzpicture}
   \tkzGetNodes
   \tkzDrawLine(A,B)
   \tkzDrawPoints(A,B)
   \tkzLabelPoints(A,B)
  \end{tikzpicture}
\end{minipage}
\end{mybox}

\subsection{要点}

以下几点对本说明文档中的大多数代码都至关重要。

\begin{itemize}
\item \verb|% !TEX TS-program = lualatex|。
必须使用 \code{lualatex} 进行编译。如果你的编辑器已据此配置，则此行可省略。

\item \verb|\usepackage[mini]{tkz-euclide}|。
\tkzNamePack{tkz-euclide} 的使用是可选的；但若加载它，建议使用 \code{mini} 选项。
该选项仅加载绘图宏，而所有计算都由 \code{lua} 处理。
\textbf{重要提示：}在目前阶段，部分绘图宏尚未完全独立于计算宏。
如果遇到编译问题，只需不加 \code{mini} 选项加载 \tkzNamePack{tkz-euclide}。

\item \verb|\usepackage{tkz-elements}|。
此宏包为必需。它提供基于 Lua 的几何引擎以及与 \tkzNamePack{tkz-euclide} 配合使用的辅助宏。

\item \verb|\directlua{...}|。
所有几何定义与计算都放在这个宏内部（或放在 \texttt{luacode} 环境中，参见 \tkzNamePack{luacode} 宏包）。

\item \verb|init_elements()|。
每个 Lua 段的开头都应调用此函数。
它会重置内部表，并清除 \tkzNamePack{tkz-elements} 所使用的全局变量。

\item \verb|\tkzGetNodes|。
此宏必须放在 \code{tikzpicture} 环境的开头。
它把在 Lua 中定义的点作为可用节点传递给 TikZ。
\end{itemize}



\subsection{测试}

要测试你的安装并跟随本文档中的示例操作，需要加载两个宏包：\tkzNamePack{tkz-euclide} 和 \tkzNamePack{tkz-elements}。第一个宏包会自动加载 \tkzNamePack{\TIKZ}，而后者对所有图形渲染都是必需的。


\tkzname{Lua} 代码作为参数传给 \tkzMacro{lualatex}{directlua} 宏。我们通常把这段代码块称为 \tkzname{Lua 部分} \footnote{这段代码也可以放在一个外部文件中，例如 \texttt{file.lua}。}。这一部分完全依赖于 \tkzNamePack{tkz-elements} 宏包。

\tkzEnv{tikz}{tikzpicture} 环境中的一个关键组件是 \tkzMacro{tkz-elements}{tkzGetNodes} 宏。该宏通过创建相应节点，把在 \tkzname{Lua} 中定义的点传递给 \tkzNamePack{\TIKZ}。所有这类点都存放在一个名为 \tkzname{z} 的表中，并使用 \tkzname{z.label} 的语法访问。这些标签随后会在 \tkzNamePack{tkz-euclide} 中重复使用。

当你通过为一个点赋予标签和坐标来定义它时，它在内部被表示为一个复数——即该点的附标。这种表示方式使得该点能够在一个正交直角坐标系中定位。

如果你想用另一种方式来渲染对象，那么需要修改的就是这个宏。例如，第~\ref{sec:metapost} 节介绍了 \tkzMacro{tkz-elements}{tkzGetNodesMP}，它是一个支持与 \code{MetaPost} 通信的变体。

另一个关键要素是使用 \tkzFct{tkz-elements}{init\_elements()} 函数，它在处理多个图形时会清除内部表\footnote{所有几何对象都存放在 Lua 表中。这些表必须定期清理，尤其是在连续创建多个图形时。}。

如果前面的代码一切正常，那么你就可以开始创建几何对象了。第~\ref{sec:class_and_object} 节介绍了可用的选项和对象结构。

最后，熟悉 \tkzNamePack{tkz-euclide} 中的基本绘图命令也很重要，因为它们将用于渲染在 \tkzname{Lua} 中定义的对象。
\endinput
